\subsection{Complementary}
\label{sec:theory_general_complementary}

There are primarily two types of transistor technologies used in
analog circuit design: NPN and PNP for Bipolar Junction Transistors
(BJTs), and NMOS and PMOS for Metal-Oxide-Semiconductor Field
Effect Transistors (MOSFETs).

These transistor types are complementary to each other, meaning
their electrical characteristics and behaviors are opposite in
nature. Any circuit designed involving transistors can be converted
into its complementary form by following these general rules:

\begin{enumerate}
  \item Replace NPN BJTs with PNP BJTs and vice versa.
  \item Replace NMOS MOSFETs with PMOS MOSFETs and vice versa.
  \item Reverse the polarity of all voltage supplies.
  \item Reverse the direction of all current sources.
\end{enumerate}

\noindent A good things to remember when designing,
\begin{itemize}
  \item N-type devices (NPN, NMOS) are typically used to
  sink conventional current to ground from a higher potential,
  while P-type devices (PNP, PMOS) are used to source conventional
  current from a higher potential toward the load.
  \item N-type devices are commonly used as low-side switches,
  while P-type devices are commonly used as high-side switches.
  This means N-type devices are often placed between the load
  and ground (or lower potential), while P-type devices are placed
  between the power supply and the load.
\end{itemize}
